\documentclass[UTF8,zihao=-4,scheme=chinese]{ctexart}

\usepackage[
  a4paper,
  left=1.55cm,
  right=1.55cm,
  top=1.65cm,
  bottom=1.8cm,
  headheight=14pt,
  headsep=0.45cm,
  footskip=0.85cm
]{geometry}
\usepackage{setspace}
\setstretch{1.12}
\setlength{\parindent}{2em}
\setlength{\parskip}{0.25em}

\usepackage{amsmath,amssymb,mathtools,bm}
\usepackage{tikz}
\usetikzlibrary{arrows.meta}
\usepackage{pgfplots}
\pgfplotsset{compat=1.18}
\usepackage{booktabs,array,tabularx}
\usepackage{enumitem}
\setlist[itemize]{leftmargin=2em,itemsep=0.15em,topsep=0.25em}
\setlist[enumerate]{leftmargin=2.2em,itemsep=0.15em,topsep=0.25em}
\usepackage[hidelinks]{hyperref}
\hypersetup{pdftitle={高等数学连续性复习讲义}, pdfauthor={Course Alchemist}}

\usepackage{fancyhdr}
\pagestyle{fancy}
\fancyhf{}
\fancyhead[L]{\small 高等数学复习讲义}
\fancyhead[R]{\small \leftmark}
\fancyfoot[C]{\small \thepage}
\renewcommand{\headrulewidth}{0.4pt}
\renewcommand{\footrulewidth}{0pt}

\usepackage[most]{tcolorbox}
\tcbuselibrary{breakable,skins,theorems}

\definecolor{CAInk}{HTML}{1F2937}
\definecolor{CASubtle}{HTML}{F8FAFC}
\definecolor{CADef}{HTML}{2563EB}
\definecolor{CAThm}{HTML}{7C3AED}
\definecolor{CALem}{HTML}{0891B2}
\definecolor{CAHint}{HTML}{059669}
\definecolor{CAExam}{HTML}{D97706}
\definecolor{CAPitfall}{HTML}{DC2626}
\definecolor{CANote}{HTML}{475569}
\definecolor{CAMethod}{HTML}{0F766E}
\definecolor{CAExample}{HTML}{4F46E5}

\tcbset{
  ca-base/.style={
    enhanced,
    breakable,
    sharp corners,
    boxrule=0.6pt,
    leftrule=2.4pt,
    arc=1.2mm,
    left=1.1mm,
    right=1.1mm,
    top=0.9mm,
    bottom=0.9mm,
    before skip=0.65em,
    after skip=0.65em,
    fonttitle=\bfseries,
    coltitle=CAInk
  }
}

\newtcbtheorem[number within=section]{definition}{定义}{ca-base,colback=CADef!5,colframe=CADef,colbacktitle=CADef!12}{def}
\newtcbtheorem[number within=section]{theorem}{定理}{ca-base,colback=CAThm!5,colframe=CAThm,colbacktitle=CAThm!12}{thm}
\newtcbtheorem[number within=section]{lemma}{引理}{ca-base,colback=CALem!5,colframe=CALem,colbacktitle=CALem!12}{lem}
\newtcbtheorem[number within=section]{hint}{提示}{ca-base,colback=CAHint!5,colframe=CAHint,colbacktitle=CAHint!12}{hint}
\newtcbtheorem[number within=section]{examspot}{考点}{ca-base,colback=CAExam!7,colframe=CAExam,colbacktitle=CAExam!14}{exam}
\newtcbtheorem[number within=section]{pitfall}{易错点}{ca-base,colback=CAPitfall!5,colframe=CAPitfall,colbacktitle=CAPitfall!12}{pit}
\newtcbtheorem[number within=section]{attention}{关注}{ca-base,colback=CANote!6,colframe=CANote,colbacktitle=CANote!14}{att}
\newtcbtheorem[number within=section]{method}{方法}{ca-base,colback=CAMethod!5,colframe=CAMethod,colbacktitle=CAMethod!12}{mtd}
\newtcbtheorem[number within=section]{example}{例题}{ca-base,colback=CAExample!5,colframe=CAExample,colbacktitle=CAExample!12}{ex}

\newtcolorbox{solutionbox}[1][]{%
  ca-base,
  title={解答},
  colback=CASubtle,
  colframe=CAInk!55,
  colbacktitle=CAInk!8,
  #1
}

\newcommand{\R}{\mathbb{R}}
\newcommand{\dd}{\,\mathrm{d}}

\title{\bfseries 高等数学：连续性复习讲义}
\author{Course Alchemist 示例项目}
\date{\today}

\begin{document}

\maketitle
\tableofcontents

\section{学习目标与位置}

连续性连接了极限、函数性质和存在性定理。复习时要同时掌握三件事：怎样用定义判定连续，怎样处理分段函数的分界点，以及怎样用连续性推出零点或不动点等存在性结论。

\begin{examspot}{本节常考方式}{continuity-exam}
连续性常与分段函数参数确定、零点存在性、方程根的存在性和反例判断结合考查。看到“闭区间连续”“端点异号”“取遍中间值”等关键词时，应优先想到介值定理。
\end{examspot}

\section{定义与判定}

\begin{definition}{连续性}{continuity}
设函数 $f$ 在点 $x_0$ 的某邻域内有定义。若
\[
  \lim_{x\to x_0} f(x)=f(x_0),
\]
则称 $f$ 在 $x_0$ 处连续。若 $f$ 在区间上每一点都连续，则称 $f$ 在该区间上连续。
\end{definition}

判定一点连续时，要同时检查三个对象：函数值 $f(x_0)$ 存在，极限 $\lim_{x\to x_0}f(x)$ 存在，二者相等。对分段函数，分界点通常是唯一需要重点检查的位置。

\begin{method}{分段函数连续性检查}{piecewise-method}
检查分段函数在分界点 $x_0$ 处连续，一般按三步进行：
\begin{enumerate}
  \item 计算左极限 $\lim_{x\to x_0^-}f(x)$；
  \item 计算右极限或函数值 $f(x_0)$；
  \item 比较左右极限与函数值是否一致。
\end{enumerate}
\end{method}

\begin{example}{确定参数}{piecewise-example}
设
\[
f(x)=
\begin{cases}
ax+1, & x<1,\\
x^2+b, & x\ge 1.
\end{cases}
\]
若 $f$ 在 $x=1$ 处连续，求 $a,b$ 满足的关系。
\end{example}

\begin{solutionbox}
由左极限与函数值相等可得
\[
  \lim_{x\to 1^-}(ax+1)=f(1).
\]
即
\[
  a+1=1+b,
\]
所以 $a=b$。
\end{solutionbox}

\section{连续函数的存在性工具}

\begin{theorem}{介值定理}{ivt}
若 $f$ 在闭区间 $[a,b]$ 上连续，且 $y$ 介于 $f(a)$ 与 $f(b)$ 之间，则存在 $\xi\in[a,b]$，使得 $f(\xi)=y$。
\end{theorem}

\begin{lemma}{零点存在性}{zero}
若 $f$ 在 $[a,b]$ 上连续，且 $f(a)f(b)<0$，则存在 $\xi\in(a,b)$，使得 $f(\xi)=0$。
\end{lemma}

\begin{hint}{构造辅助函数}{auxiliary}
存在性证明常把目标方程改写为 $g(x)=0$。先构造连续函数 $g$，再检查端点符号或中间值条件。
\end{hint}

\begin{figure}[htbp]
\centering
\begin{tikzpicture}[scale=0.95]
  \draw[-{Latex}] (-0.2,0) -- (5.2,0) node[right] {$x$};
  \draw[-{Latex}] (0,-1.4) -- (0,2.2) node[above] {$y$};
  \draw[thick,CAExample,domain=0.4:4.6,smooth,samples=80] plot (\x,{0.25*(\x-1.2)*(\x-4)+0.2});
  \draw[dashed] (0.4,0) node[below] {$a$} -- (0.4,{0.25*(0.4-1.2)*(0.4-4)+0.2});
  \draw[dashed] (4.6,0) node[below] {$b$} -- (4.6,{0.25*(4.6-1.2)*(4.6-4)+0.2});
  \draw[dashed] (2.15,0) node[below] {$\xi$} -- (2.15,{0.25*(2.15-1.2)*(2.15-4)+0.2});
  \node[left] at (0,{0.25*(0.4-1.2)*(0.4-4)+0.2}) {$f(a)$};
  \node[left] at (0,{0.25*(4.6-1.2)*(4.6-4)+0.2}) {$f(b)$};
\end{tikzpicture}
\caption{连续曲线从负值到正值时必须穿过零点。}
\end{figure}

\begin{pitfall}{忽略闭区间连续性}{ivt-pitfall}
介值定理要求函数在整个闭区间上连续。只知道端点函数值异号，但函数中间有断点，不能直接推出零点存在。
\end{pitfall}

\section{分层练习}

\begin{enumerate}
  \item[1.]（基础掌握，6 分）设
  \[
    f(x)=
    \begin{cases}
      ax+1, & x<1,\\
      x^2+b, & x\ge 1.
    \end{cases}
  \]
  若 $f$ 在 $x=1$ 处连续，求 $a,b$ 满足的关系。

  \item[2.]（考试常规，8 分）设 $f$ 在 $[0,1]$ 上连续，且 $f(0)<0<f(1)$。证明存在 $\xi\in(0,1)$，使得 $f(\xi)=0$。

  \item[3.]（拔高区分，10 分）设 $f$ 在 $[0,1]$ 上连续，且对任意 $x\in[0,1]$ 有 $0\le f(x)\le 1$。证明存在 $\xi\in[0,1]$，使得 $f(\xi)=\xi$。
\end{enumerate}

\section{练习解析}

\begin{solutionbox}[title={第 1 题}]
连续要求分界点处左极限与函数值相等。由
\[
  \lim_{x\to 1^-}(ax+1)=a+1,\qquad f(1)=1+b
\]
可得 $a+1=1+b$，所以 $a=b$。
\end{solutionbox}

\begin{solutionbox}[title={第 2 题}]
因为 $f$ 在闭区间 $[0,1]$ 上连续，且 $f(0)<0<f(1)$，所以 $0$ 介于 $f(0)$ 与 $f(1)$ 之间。由介值定理，存在 $\xi\in(0,1)$，使得 $f(\xi)=0$。
\end{solutionbox}

\begin{solutionbox}[title={第 3 题}]
令 $g(x)=f(x)-x$。由于 $f$ 连续，$g$ 在 $[0,1]$ 上连续。又因为 $0\le f(0)\le 1$，所以
\[
  g(0)=f(0)\ge 0.
\]
同时 $0\le f(1)\le 1$，所以
\[
  g(1)=f(1)-1\le 0.
\]
若端点已有 $g(0)=0$ 或 $g(1)=0$，结论成立；否则由介值定理，存在 $\xi\in(0,1)$ 使 $g(\xi)=0$，即 $f(\xi)=\xi$。
\end{solutionbox}


\end{document}
